\documentclass[9pt,a4paper]{extarticle}

\usepackage[T1]{fontenc}
\usepackage[utf8]{inputenc}
\usepackage{babel}
\usepackage[margin=1.2cm,top=0.8cm,bottom=0.7cm]{geometry}
\usepackage{enumitem}
\usepackage[hidelinks]{hyperref}
\usepackage{titlesec}

\pagestyle{empty}
\setlength{\parindent}{0pt}
\setlength{\parskip}{0pt}
\renewcommand{\familydefault}{\sfdefault}

\titleformat{\section}{\normalsize\bfseries\scshape}{}{0pt}{}[\vspace{-0.6em}\rule{\linewidth}{0.4pt}]
\titlespacing*{\section}{0pt}{0.4em}{0.25em}

\setlist[itemize]{leftmargin=1.1em,itemsep=0em,topsep=0.1em,parsep=0pt,partopsep=0pt}

% \entry{Title}{Right-hand text}{Subtitle}{Right-hand subtitle}
\newcommand{\entry}[4]{%
  \textbf{#1}\hfill #2\\[-0.1em]
  \textit{#3}\hfill \textit{#4}\par\vspace{0.1em}}

\begin{document}

\begin{center}
  {\LARGE\bfseries Charly Alizadeh}\\[0.3em]
  Ingénieur automatisation réseau \;|\; Région parisienne, France\\[0.2em]
  \href{mailto:alizadehcharly@gmail.com}{alizadehcharly@gmail.com} \;|\;
  \href{https://www.linkedin.com/in/charly-alizadeh}{linkedin.com/in/charly-alizadeh} \;|\;
  \href{https://github.com/spikte}{github.com/spikte} \;|\;
  \href{https://github.com/charlyalizadeh}{github.com/charlyalizadeh}
\end{center}

\section{Expérience professionnelle}

\entry{Ingénieur Systèmes et Réseaux}{mai 2025 -- août 2025}{Axeria}{Lyon, France}
\begin{itemize}
  \item Conception et mise en œuvre de scripts d'automatisation s'appuyant sur les API Microsoft Graph et Tenable
        (PowerShell, Python) afin d'industrialiser des tâches manuelles~: correctifs, enrôlement Autopilot,
        boîtes aux lettres, épuration des postes inactifs.
  \item Développement d'une application web interne (Python/Flask) agrégeant les données issues d'Intune,
        d'Active Directory, d'Entra ID, de Tenable et de ManageEngine pour une vision consolidée du parc.
  \item Administration des identités (Active Directory, Entra ID)~; rédaction de procédures et de standards IT.
\end{itemize}


\entry{Ingénieur de recherche}{mars 2024 -- sept. 2024}{Université Sorbonne Paris Nord}{Paris, France}
\begin{itemize}
  \item Implémentation en Julia d'un modèle d'Optimal Power Flow (OPF) et de ses relaxations semi-définies (SDP).
  \item Création d'une base de données de décompositions en cliques et apprentissage automatique pour sélectionner la plus
        favorable dans le but d'accélérer la résolution des relaxations SDP.
\end{itemize}


\entry{Chargé de travaux dirigés}{janv. 2023 -- janv. 2024}{ESILV -- École Supérieure d'Ingénieurs Léonard de Vinci}{Courbevoie, France}
\begin{itemize}
  \item Encadrement de TD/TP de la L1 au M1~: algorithmique et graphes, C\#, machine learning, deep learning,
        science des données~; suivi de projets et évaluation des travaux.
\end{itemize}

\section{Projets personnels}

\entry{Fabric de data center «~as code~»}{Go, Ansible, Jinja2, containerlab, FRRouting, Docker}{}{\href{https://github.com/spikte/data_center_fabric}{github.com/spikte/data\_center\_fabric}}
\begin{itemize}
  \item Contrôleur Go générant, à partir d'une description YAML déclarative de la fabric (source de vérité unique),
        la topologie containerlab et l'inventaire Ansible correspondant.
  \item Rôles Ansible produisant les configurations FRRouting par templates Jinja2 (underlay et overlay eBGP).
  \item Fabric leaf-spine multi-tenant, overlay VXLAN EVPN (L2VNI et L3VNI)~; serveurs émulés par des bridges
        Linux VLAN-aware et des conteneurs netshoot~; validation automatisée de la connectivité de bout en bout.
\end{itemize}


\entry{Netsim -- Simulateur réseau}{C++, Qt}{}{\href{https://github.com/spikte/netsim}{github.com/spikte/netsim}}
\begin{itemize}
  \item Simulateur entièrement développé sans bibliothèque réseau~: commutation L2 et table MAC, ARP, VLAN,
        LLDP et STP partiels, routage IPv4, ICMP.
  \item CLI inspirée de Cisco IOS (\texttt{show ip route}, \texttt{show vlan}, routes statiques, \texttt{ping},
        \texttt{traceroute}) pour observer et diagnostiquer le réseau simulé.
  \item Interface graphique en Qt
\end{itemize}


\entry{packet\_viz -- Analyse de l'origine du trafic}{C, raw sockets, MaxMind GeoLite2}{}{\href{https://github.com/charlyalizadeh/packet_viz}{github.com/charlyalizadeh/packet\_viz}}
\begin{itemize}
  \item Capture du trafic TCP/IPv4 par sockets bruts~; identification du système autonome (AS) et de la localisation géographique de chaque source par recherche
        du plus long préfixe correspondant (longest-prefix match) dans un trie IP construit à partir des tables MaxMind.
\end{itemize}


\entry{spikte\_chat -- Serveur de messagerie}{C++, sockets Linux, epoll, TLS, SQLite}{}{\href{https://github.com/spikte/spikte_chat}{github.com/spikte/spikte\_chat}}
\begin{itemize}
  \item Serveur multi-clients fondé sur une boucle d'événements epoll (sans thread par client)~; transport TCP ou TLS,
        salons multiples, persistance SQLite, mots de passe hachés avec Argon2.
\end{itemize}

\section{Formation}
\entry{Systèmes informatiques et réseaux}{sept. 2025 -- juil. 2026}{Conservatoire National des Arts et Métiers (CNAM)}{Paris, France}
\entry{Diplôme d'ingénieur, parcours recherche -- Big Data et intelligence artificielle}{2017 -- 2022}{ESILV -- École Supérieure d'Ingénieurs Léonard de Vinci}{Courbevoie, France}
\entry{Mathématiques et informatique -- échange universitaire}{2019 -- 2020}{Chung-Ang University}{Séoul, Corée du Sud}

\section{Compétences techniques}
\begin{tabular}{@{}p{3.3cm}p{13.7cm}@{}}
\textbf{Automatisation} & Ansible, Jinja2, Go, Python, PowerShell, YAML déclaratif \\
\textbf{Routage et data center} & eBGP, VXLAN EVPN (L2VNI, L3VNI), architecture leaf-spine, multi-tenant, IPv4, FRRouting, SRLinux \\
\textbf{Commutation} & Ethernet, VLAN, bridges Linux VLAN-aware, MAC/ARP, LLDP, STP \\
\textbf{Labs et outillage} & containerlab, Docker, sockets Linux, epoll, TLS, Git, CMake \\
\textbf{Systèmes et sécurité} & Linux, Windows, Active Directory, Entra ID, Intune, Autopilot, Tenable, ManageEngine \\
\textbf{Autres langages} & C, C++, Julia, SQL, C\# \\
\end{tabular}

\vspace{0.3em}
{\small\textbf{Publication~:} C.~Alizadeh et al., \textit{How to Learn the Optimal Clique Decompositions in Solving Semidefinite Relaxations for OPF},
IJCNN 2022.
\href{https://ieeexplore.ieee.org/abstract/document/9891912}{DOI~: 10.1109/IJCNN55064.2022.9891912}}

\end{document}
